\section{Тест производительности}

В данном разделе проводится анализ эффективности реализованной структуры данных PATRICIA 
в сравнении со стандартным ассоциативным контейнером \texttt{std::map}.

Тестирование выполнялось на случайно сгенерированных данных. Ключи представляют собой строки 
длиной до 256 символов, состоящие из букв английского алфавита.

Замерялось суммарное время выполнения операций вставки, поиска и удаления.

\begin{alltt}
--------------------------------------------------------
| Elements   | Patricia           | std::map           |
--------------------------------------------------------
|     100000 |             107 ms |             138 ms |
|     500000 |            1034 ms |            1194 ms |
|    1000000 |            1720 ms |            2715 ms |
--------------------------------------------------------
\end{alltt}

Как видно из результатов, структура PATRICIA демонстрирует более высокую производительность 
по сравнению с \texttt{std::map}, особенно на больших объёмах данных.

Это объясняется следующими причинами:

\begin{itemize}
    \item Используется побитовое сравнение ключей вместо посимвольного;
    
    \item Отсутствует зависимость от числа элементов ($\log N$), как в сбалансированных деревьях;
    
    \item Структура эффективно разделяет ключи по первому различающемуся биту.
\end{itemize}

Таким образом, bitwise PATRICIA является эффективной структурой данных для хранения 
и поиска строковых ключей.